% GENERATED FILE. Edit canonical lessons, then run npm run generate:books.
% canonical-source-hash: fnv1a64:e98957998b7e7442
% canonical-lessons: RU-C451-posolstvo, RU-C451-zagranpasport, RU-C451-banderol, RU-C451-poluchatel, RU-C451-vydavat

\chapter{The embassy and the post}
\label{ch:ru-the-embassy-and-the-post}
\hlchaptermodality{451}

\begin{chapteropening}
\textbf{By the end of this chapter:} \emph{I can visit the embassy and send small parcels: the embassy, a passport for travel abroad, a small parcel, the addressee and issuing documents.}

Complete the last lesson of chapter 451: I can visit the embassy and send small parcels: the embassy, a passport for travel abroad, a small parcel, the addressee and issuing documents.
\end{chapteropening}

\section[\texorpdfstring{\ru{посольство} (posólstvo)}{посольство (posólstvo)}]{\ru{посольство} (posólstvo) — an embassy}

\label{lesson:RU-C451-posolstvo}



\begin{quote}
Before the new one: say the Russian for to get, then the Russian for to certify.
\end{quote}

\subsection*{You'll want to know: \ru{посольство}}
\textbf{\ru{посольство}} (\emph{posólstvo}) — \textquotedblleft{}an embassy\textquotedblright{}.

\begin{quote}
\ru{Мы} \ru{пошли} \ru{в} \ru{посольство} \ru{за} \ru{визой}.
\end{quote}

\begin{quote}
\emph{My poshlí v posólstvo za vízoy.}
\end{quote}

\begin{quote}
We went to the embassy for a visa.
\end{quote}

Neuter. Built on \emph{\ru{посол}} — an ambassador.

\begin{cousinweb}
\emph{\ru{посольство}} is built on \emph{\ru{посол}} 'an envoy, an ambassador' with \emph{-\ru{ство}}. \emph{\ru{Посол}} goes with \emph{\ru{посылать}} 'to send': one who is sent. EtymDB lists \emph{\ru{посольство}} without a deeper source.
\end{cousinweb}

\begin{grammarlens}[title={What you've built}]
One more word for telling, explaining and describing at B1.
\end{grammarlens}

\subsection*{Your turn}
\begin{itemize}
\raggedright
  \item \emph{Say it:} \emph{posólstvo}
  \item \emph{Say it:} \emph{posólstvo}, once more
  \item \emph{Say it:} a sentence of your own with \emph{posólstvo}
  \item \emph{Recall:} say the Russian for to get, then the Russian for to certify, then say \emph{posólstvo} again
\end{itemize}

\begin{tcolorbox}[breakable,colback=teal!4,colframe=teal!35!black,title={Before you move on}]
What does \ru{посольство} mean? (\textquotedblleft{}An embassy\textquotedblright{}.) Say it once more.
\end{tcolorbox}
\section[\texorpdfstring{\ru{загранпаспорт} (zagranpásport)}{загранпаспорт (zagranpásport)}]{\ru{загранпаспорт} (zagranpásport) — a passport for travel abroad}

\label{lesson:RU-C451-zagranpasport}



\begin{quote}
Before the new one: say the Russian for to certify, then the Russian for an embassy.
\end{quote}

\subsection*{You'll want to know: \ru{загранпаспорт}}
\textbf{\ru{загранпаспорт}} (\emph{zagranpásport}) — \textquotedblleft{}a passport for travel abroad\textquotedblright{}.

\begin{quote}
\ru{Для} \ru{поездки} \ru{за} \ru{границу} \ru{нужен} \ru{загранпаспорт}.
\end{quote}

\begin{quote}
\emph{Dlya poyézdki za granítsu núzhen zagranpásport.}
\end{quote}

\begin{quote}
You need a passport to travel abroad.
\end{quote}

Short for \emph{\ru{заграничный} \ru{паспорт}}; Russians also have an internal passport.

\begin{cousinweb}
\emph{\ru{загранпаспорт}} is a short form of \emph{\ru{заграничный} \ru{паспорт}}, from \emph{\ru{за} \ru{границей}} 'abroad', on \emph{\ru{граница}} 'a border', and \emph{\ru{паспорт}} 'a passport'.
\end{cousinweb}

\begin{grammarlens}[title={What you've built}]
One more word for telling, explaining and describing at B1.
\end{grammarlens}

\subsection*{Your turn}
\begin{itemize}
\raggedright
  \item \emph{Say it:} \emph{zagranpásport}
  \item \emph{Say it:} \emph{zagranpásport}, once more
  \item \emph{Say it:} a sentence of your own with \emph{zagranpásport}
  \item \emph{Recall:} say the Russian for to certify, then the Russian for an embassy, then say \emph{zagranpásport} again
\end{itemize}

\begin{tcolorbox}[breakable,colback=teal!4,colframe=teal!35!black,title={Before you move on}]
What does \ru{загранпаспорт} mean? (\textquotedblleft{}A passport for travel abroad\textquotedblright{}.) Say it once more.
\end{tcolorbox}
\section[\texorpdfstring{\ru{бандероль} (banderól)}{бандероль (banderól)}]{\ru{бандероль} (banderól) — a small parcel; printed matter (by post)}

\label{lesson:RU-C451-banderol}



\begin{quote}
Before the new one: say the Russian for an embassy, then the Russian for a passport for travel abroad.
\end{quote}

\subsection*{You'll want to know: \ru{бандероль}}
\textbf{\ru{бандероль}} (\emph{banderól}) — \textquotedblleft{}a small parcel; printed matter (by post)\textquotedblright{}.

\begin{quote}
\ru{Я} \ru{отправил} \ru{книги} \ru{бандеролью}.
\end{quote}

\begin{quote}
\emph{Ya otprávil knígi banderólyu.}
\end{quote}

\begin{quote}
I sent the books as a small parcel.
\end{quote}

Feminine. Smaller and cheaper than a \emph{\ru{посылка}}.

\begin{cousinweb}
\emph{\ru{бандероль}} was borrowed from French \emph{banderole}, from Italian \emph{banderuola} (EtymDB), a little banner; the word probably first named the paper band round a packet.
\end{cousinweb}

\begin{grammarlens}[title={What you've built}]
One more word for telling, explaining and describing at B1.
\end{grammarlens}

\subsection*{Your turn}
\begin{itemize}
\raggedright
  \item \emph{Say it:} \emph{banderól}
  \item \emph{Say it:} \emph{banderól}, once more
  \item \emph{Say it:} a sentence of your own with \emph{banderól}
  \item \emph{Recall:} say the Russian for an embassy, then the Russian for a passport for travel abroad, then say \emph{banderól} again
\end{itemize}

\begin{tcolorbox}[breakable,colback=teal!4,colframe=teal!35!black,title={Before you move on}]
What does \ru{бандероль} mean? (\textquotedblleft{}A small parcel; printed matter (by post)\textquotedblright{}.) Say it once more.
\end{tcolorbox}
\section[\texorpdfstring{\ru{получатель} (poluchátel)}{получатель (poluchátel)}]{\ru{получатель} (poluchátel) — an addressee, a recipient}

\label{lesson:RU-C451-poluchatel}



\begin{quote}
Before the new one: say the Russian for a passport for travel abroad, then the Russian for a small parcel.
\end{quote}

\subsection*{You'll want to know: \ru{получатель}}
\textbf{\ru{получатель}} (\emph{poluchátel}) — \textquotedblleft{}an addressee, a recipient\textquotedblright{}.

\begin{quote}
\ru{Посылка} \ru{не} \ru{дошла} \ru{до} \ru{получателя}.
\end{quote}

\begin{quote}
\emph{Posýlka ne doshlá do poluchátelya.}
\end{quote}

\begin{quote}
The parcel did not reach the addressee.
\end{quote}

Masculine. The sender is the \emph{\ru{отправитель}}.

\begin{cousinweb}
\emph{\ru{получатель}} is built on \emph{\ru{получать}} 'to receive' with \emph{-\ru{тель}} 'one who'. EtymDB lists \emph{\ru{получатель}} without a deeper source.
\end{cousinweb}

\begin{grammarlens}[title={What you've built}]
One more word for telling, explaining and describing at B1.
\end{grammarlens}

\subsection*{Your turn}
\begin{itemize}
\raggedright
  \item \emph{Say it:} \emph{poluchátel}
  \item \emph{Say it:} \emph{poluchátel}, once more
  \item \emph{Say it:} a sentence of your own with \emph{poluchátel}
  \item \emph{Recall:} say the Russian for a passport for travel abroad, then the Russian for a small parcel, then say \emph{poluchátel} again
\end{itemize}

\begin{tcolorbox}[breakable,colback=teal!4,colframe=teal!35!black,title={Before you move on}]
What does \ru{получатель} mean? (\textquotedblleft{}An addressee, a recipient\textquotedblright{}.) Say it once more.
\end{tcolorbox}
\section[\texorpdfstring{\ru{выдавать} (vydavát)}{выдавать (vydavát)}]{\ru{выдавать} (vydavát) — to issue, to hand out}

\label{lesson:RU-C451-vydavat}



\begin{quote}
Before the new one: say the Russian for a small parcel, then the Russian for an addressee.
\end{quote}

\subsection*{You'll want to know: \ru{выдавать}}
\textbf{\ru{выдавать}} (\emph{vydavát}) — \textquotedblleft{}to issue, to hand out\textquotedblright{}.

\begin{quote}
\ru{Посольство} \ru{выдаёт} \ru{визы} \ru{за} \ru{неделю}.
\end{quote}

\begin{quote}
\emph{Posólstvo vydayót vízy za nedélyu.}
\end{quote}

\begin{quote}
The embassy issues visas within a week.
\end{quote}

Imperfective; perfective \emph{\ru{выдать}}. Present \emph{\ru{выдаю}, \ru{выдаёшь}}.

\begin{cousinweb}
\emph{\ru{выдавать}} is built from \emph{\ru{вы}-} 'out' and \emph{\ru{давать}} 'to give', which is inherited from Proto-Slavic \emph{dāvati} (EtymDB).
\end{cousinweb}

\begin{grammarlens}[title={What you've built}]
One more word for telling, explaining and describing at B1.
\end{grammarlens}

\subsection*{Your turn}
\begin{itemize}
\raggedright
  \item \emph{Say it:} \emph{vydavát}
  \item \emph{Say it:} \emph{vydavát}, once more
  \item \emph{Say it:} a sentence of your own with \emph{vydavát}
  \item \emph{Recall:} say the Russian for a small parcel, then the Russian for an addressee, then say \emph{vydavát} again
\end{itemize}

\begin{tcolorbox}[breakable,colback=teal!4,colframe=teal!35!black,title={Before you move on}]
What does \ru{выдавать} mean? (\textquotedblleft{}To issue, to hand out\textquotedblright{}.) Say it once more.
\end{tcolorbox}
